
\subsubsection{6.1 Interpretation of Results}

\textbf{Inductive bias differences are quantifiable.} DWS and NFT produce measurably different weight update patterns (CoV 1.44 vs 1.35). This provides the first concrete operationalization of locality versus global attention in weight-space learning.

\textbf{NFT global attention benefits accuracy prediction.} NFT achieves 12.3\% lower RMSE than DWS on the synthetic accuracy task. The target was computed from global weight statistics (mean norms and means across layers), which aligns with NFT's architectural strength in capturing aggregate information.

\textbf{DWS locality advantage on backdoor detection is unverified.} The synthetic backdoor signal was not learnable by any architecture. This is an experimental design limitation, not evidence against the hypothesis. The localized perturbations were likely absorbed by normalization or did not match patterns that the architectures could learn.

\textbf{The interaction effect exists but is partial.} The ANOVA indicates a significant architecture-task interaction (p<0.001), but only one direction of the predicted crossover pattern was observed.

\subsubsection{6.2 Limitations}

\begin{enumerate}
\item \textbf{Synthetic backdoor signals were unlearnable.} The localized row perturbation approach does not produce signals that weight-space architectures can detect. Future work should use real TrojAI benchmark data with genuine backdoor signatures.
\end{enumerate}

\begin{enumerate}
\item \textbf{Dataset ceiling effects.} The MNIST-INR classification task achieves 100\% accuracy for all architectures, preventing differentiation on that dimension and invalidating sample efficiency tests.
\end{enumerate}

\begin{enumerate}
\item \textbf{Parameter count mismatch.} DWS (4.9M), NFT (5.3M), and MLP (9.8M) are not matched within 10\%. However, MLP with approximately twice the parameters underperforms NFT on accuracy prediction, suggesting inductive bias rather than capacity drives the observed difference.
\end{enumerate}

\begin{enumerate}
\item \textbf{Synthetic targets.} The accuracy prediction target is computed from weight statistics rather than actual model performance, which may not reflect real accuracy prediction difficulty.
\end{enumerate}

\subsubsection{6.3 Implications}

For practitioners selecting weight-space architectures:

\begin{itemize}
\item \textbf{Accuracy prediction and aggregate statistics:} NFT or similar global attention architectures are preferable based on the observed 12.3\% RMSE improvement.
\end{itemize}

\begin{itemize}
\item \textbf{Backdoor detection and local anomaly tasks:} No recommendation can be made from this study. Evaluation on real TrojAI benchmarks is required.
\end{itemize}

The methodology for comparing inductive biases via training dynamics (CoV analysis) may be applicable beyond weight-space learning.
